% ================================================================
\section{Introduction: A New Kind of Mathematical Object}
\label{sec:introduction}
% ================================================================

In the taxonomy of mathematical claims, a proposition is traditionally
assigned one of three epistemic statuses: \emph{proved} (its truth is
established by a chain of reasoning from accepted axioms),
\emph{refuted} (its negation is proved), or \emph{open} (neither its
truth nor its falsity has been established). G\"odel's incompleteness
theorems~\cite{godel1931} introduced a fourth category---\emph{independent}---for
statements that are neither provable nor refutable within a given
formal system. The Cathedral suggests a fifth:

\begin{quote}
\emph{Formally reduced}: a claim whose logical structure has been
completely mapped by machine verification, so that its truth or
falsity depends on exactly one precisely characterized gap,
with everything else certified by a proof assistant.
\end{quote}

The Riemann Hypothesis (RH), arguably the deepest open problem in
mathematics, has been formally reduced in this sense. The Cathedral
is a Lean~4 formalization comprising 330 active files and
approximately 97{,}000 lines of code, with zero compilation errors,
that establishes the biconditional
\[
  \RH \;\iff\; d_N^2 \to 0
\]
where $d_N^2 = \inf_v \int_0^1 (1 - \sum_{k=2}^{N} v_k \{1/(kx)\})^2\,dx$
is the Nyman--Beurling distance. The converse direction
($d_N^2 \to 0 \implies \RH$) is proved with \textbf{zero} custom
axioms and \textbf{zero sorry}. The forward direction
($\RH \implies d_N^2 \to 0$) depends on exactly \textbf{one}
literature axiom~\cite{baezduarte2003}.

This paper examines what this formal reduction \emph{means}---not
mathematically (the companion paper~\cite{cathedral} provides the
technical details) nor physically (the physics dictionary~\cite{cathedral-physics}
catalogues 160+ structural correspondences between the proof and
quantum field theory), but \emph{philosophically}. We engage with
the questions that the Cathedral raises for the philosophy of
mathematics:

\begin{enumerate}
  \item \textbf{Epistemology.} What kind of knowledge is a
    machine-verified conditional equivalence? If the compiler has
    certified that RH is equivalent to a concrete analytic inequality,
    do we ``know'' more about RH than we did before---even though
    the hypothesis itself remains unproved? (\S\ref{sec:epistemology})

  \item \textbf{Ontology.} The physics dictionary reveals that every
    mathematical structure in the proof has a precise physical dual:
    the Gram matrix is a quantum correlator, the Mertens function is
    a magnetization, the Parseval Bridge is unitarity. Are these
    correspondences structural or metaphorical? What do they tell us
    about the nature of mathematical objects?
    (\S\ref{sec:ontology})

  \item \textbf{Foundations.} The Cathedral reduces an open problem
    to one axiom, but the difficulty of proving that axiom appears
    to be a \emph{topological invariant}: it cannot be eliminated
    by changing proof frameworks, only redistributed. What does this
    ``conservation of difficulty'' say about the limits of formal
    methods? (\S\ref{sec:hilbert-godel})

  \item \textbf{Cognition.} The Cathedral was built by a tripartite
    collaboration: a human mathematician, two AI systems (Claude
    and Gemini), and a proof compiler (Lean~4). No single
    participant comprehends the entire proof. What model of
    mathematical understanding does this require?
    (\S\ref{sec:cognition})

  \item \textbf{Aesthetics.} The Cathedral's naming conventions,
    architectural metaphors, and the closing poem of its physics
    paper (``the integers are the particles; the primes are the
    interactions'') treat beauty as a structural feature of proof.
    Is this legitimate, or merely decoration?
    (\S\ref{sec:aesthetics})
\end{enumerate}

We argue for a position we call \emph{moderate structural realism}:
the integers possess genuine structure that is revealed, not
constructed, by formalization; the physics correspondences are
evidence that this structure is the same structure that governs
the physical world; and the machine-verified reduction is a new and
philosophically significant form of mathematical knowledge, distinct
from both proof and conjecture.

\subsection{Scope and Method}

This paper is addressed to philosophers of mathematics in the
tradition of Benacerraf~\cite{benacerraf1973},
Shapiro~\cite{shapiro1997}, Lakatos~\cite{lakatos1976}, and
Corfield~\cite{corfield2003}---those who take the practice of
mathematics seriously as a source of philosophical insight, rather
than treating it as a mere application of formal logic. We follow
Mancosu's~\cite{mancosu2008} program of drawing philosophical
conclusions from mathematical \emph{practice}, and we take seriously
Thurston's~\cite{thurston1994} observation that mathematical
understanding is not reducible to formal proof.

At the same time, we are writing about a \emph{formal proof}---a
mathematical object that exists precisely as a term in a type theory.
This creates a productive tension: the Cathedral is simultaneously
the kind of mathematics that philosophy of mathematics should attend
to (a complex, multi-authored, computationally intensive research
program with deep connections to physics) and the kind that
traditional philosophy of mathematics can most easily analyze
(a formal derivation from axioms to conclusions, checkable by
algorithm).

We quote specific Lean declarations (e.g., \lean{gram\_bound\_implies\_rh})
and specific numerical results (e.g., $v^TGv = 0.7367$ at $N = 55{,}440$)
throughout, not as ornament but as evidence. The philosophical claims
we make are grounded in the concrete mathematical structures of the
Cathedral, and the reader is invited to verify them against the
public repository.
